\section{Morphisms into Proj}
\label{section-morphisms-proj}

\noindent
Let $S$ be a graded ring.
Let $X = \text{Proj}(S)$ be the homogeneous spectrum of $S$.
Let $d \geq 1$ be an integer.
Consider the open subscheme
\begin{equation}
\label{equation-Ud}
U_d = \bigcup\nolimits_{f  \in S_d} D_{+}(f)
\quad\subset\quad
X = \text{Proj}(S)
\end{equation}
Note that $d | d' \Rightarrow U_d \subset U_{d'}$ and
$X = \bigcup_d U_d$. Neither $X$ nor $U_d$ need
be quasi-compact, see Algebra, Lemma \ref{algebra-lemma-topology-proj}.
Let us write $\mathcal{O}_{U_d}(n) = \mathcal{O}_X(n)|_{U_d}$.
By Lemma \ref{lemma-when-invertible}
we know that $\mathcal{O}_{U_d}(nd)$, $n \in \mathbf{Z}$
is an invertible $\mathcal{O}_{U_d}$-module and
that all the multiplication maps
$\mathcal{O}_{U_d}(nd) \otimes_{\mathcal{O}_{U_d}} \mathcal{O}_{U_d}(m)
\to \mathcal{O}_{U_d}(nd + m)$ of
(\ref{equation-multiply}) are isomorphisms. In particular we have
$\mathcal{O}_{U_d}(nd) \cong \mathcal{O}_{U_d}(d)^{\otimes n}$.
The graded ring map (\ref{equation-global-sections}) on global sections
combined with restriction to $U_d$ give a homomorphism of graded rings
\begin{equation}
\label{equation-psi-d}
\psi^d : S^{(d)} \longrightarrow \Gamma_*(U_d, \mathcal{O}_{U_d}(d)).
\end{equation}
For the notation $S^{(d)}$, see Algebra, Section \ref{algebra-section-graded}.
For the notation $\Gamma_*$ see
Modules, Definition \ref{modules-definition-gamma-star}.
Moreover, since $U_d$ is covered by the opens $D_{+}(f)$, $f \in S_d$
we see that $\mathcal{O}_{U_d}(d)$ is globally generated
by the sections in the image of
$\psi^d_1 : S^{(d)}_1 = S_d \to \Gamma(U_d, \mathcal{O}_{U_d}(d))$, see
Modules, Definition \ref{modules-definition-globally-generated}.

\medskip\noindent
Let $Y$ be a scheme, and let $\varphi : Y \to X$ be a morphism of schemes.
Assume the image $\varphi(Y)$ is contained in the open subscheme
$U_d$ of $X$.
By the discussion following
Modules, Definition \ref{modules-definition-gamma-star}
we obtain a homomorphism of graded rings
$$
\Gamma_*(U_d, \mathcal{O}_{U_d}(d))
\longrightarrow
\Gamma_*(Y, \varphi^*\mathcal{O}_X(d)).
$$
The composition of this and $\psi^d$ gives a graded ring
homomorphism
\begin{equation}
\label{equation-psi-phi-d}
\psi_\varphi^d :
S^{(d)}
\longrightarrow
\Gamma_*(Y, \varphi^*\mathcal{O}_X(d))
\end{equation}
which has the property that the invertible sheaf
$\varphi^*\mathcal{O}_X(d)$ is globally generated
by the sections in the image of
$(S^{(d)})_1 = S_d \to \Gamma(Y, \varphi^*\mathcal{O}_X(d))$.

\begin{lemma}
\label{lemma-converse-construction}
Let $S$ be a graded ring, and $X = \text{Proj}(S)$.
Let $d \geq 1$ and $U_d \subset X$ as above.
Let $Y$ be a scheme.
Let $\mathcal{L}$ be an invertible sheaf on $Y$.
Let $\psi : S^{(d)} \to \Gamma_*(Y, \mathcal{L})$ be
a graded ring homomorphism such that $\mathcal{L}$ is
generated by the sections in the image of
$\psi|_{S_d} : S_d \to \Gamma(Y, \mathcal{L})$.
Then there exist a morphism
$\varphi : Y \to X$ such that $\varphi(Y) \subset U_d$ and
an isomorphism $\alpha : \varphi^*\mathcal{O}_{U_d}(d) \to \mathcal{L}$
such that $\psi_\varphi^d$ agrees with $\psi$ via $\alpha$:
$$
\xymatrix{
\Gamma_*(Y, \mathcal{L}) &
\Gamma_*(Y, \varphi^*\mathcal{O}_{U_d}(d)) \ar[l]^-\alpha &
\Gamma_*(U_d, \mathcal{O}_{U_d}(d)) \ar[l]^-{\varphi^*} \\
S^{(d)} \ar[u]^\psi & &
S^{(d)} \ar[u]^{\psi^d} \ar[ul]^{\psi^d_\varphi} \ar[ll]_{\text{id}}
}
$$
commutes. Moreover, the pair $(\varphi, \alpha)$ is unique.
\end{lemma}

\begin{proof}
Pick $f \in S_d$. Denote $s = \psi(f) \in \Gamma(Y, \mathcal{L})$.
On the open set $Y_s$ where $s$ does not vanish multiplication
by $s$ induces an isomorphism $\mathcal{O}_{Y_s} \to \mathcal{L}|_{Y_s}$,
see Modules, Lemma \ref{modules-lemma-s-open}. We will denote
the inverse of this map $x \mapsto x/s$, and similarly for
powers of $\mathcal{L}$. Using this we
define a ring map $\psi_{(f)} : S_{(f)} \to \Gamma(Y_s, \mathcal{O}_Y)$
by mapping the fraction $a/f^n$ to $\psi(a)/s^n$.
By Schemes, Lemma \ref{schemes-lemma-morphism-into-affine}
this corresponds to a morphism
$\varphi_f : Y_s \to \Spec(S_{(f)}) = D_{+}(f)$.
We also introduce the isomorphism
$\alpha_f : \varphi_f^*\mathcal{O}_{D_{+}(f)}(d) \to \mathcal{L}|_{Y_s}$
which maps the pullback of the trivializing section
$f$ over $D_{+}(f)$ to the trivializing section $s$ over $Y_s$.
With this choice the commutativity of the diagram in the lemma
holds with $Y$ replaced by $Y_s$, $\varphi$ replaced by $\varphi_f$,
and $\alpha$ replaced by $\alpha_f$; verification omitted.

\medskip\noindent
Suppose that $f' \in S_d$ is a second element, and denote
$s' = \psi(f') \in \Gamma(Y, \mathcal{L})$. Then
$Y_s \cap Y_{s'} = Y_{ss'}$ and similarly
$D_{+}(f) \cap D_{+}(f') = D_{+}(ff')$.
In Lemma \ref{lemma-ample-on-proj} we saw that
$D_{+}(f') \cap D_{+}(f)$ is the same as the set
of points of $D_{+}(f)$ where the section of
$\mathcal{O}_X(d)$ defined by $f'$ does not vanish.
Hence
$\varphi_f^{-1}(D_{+}(f') \cap D_{+}(f)) = Y_s \cap Y_{s'}
= \varphi_{f'}^{-1}(D_{+}(f') \cap D_{+}(f))$.
On $D_{+}(f) \cap D_{+}(f')$ the fraction $f/f'$ is an
invertible section of the structure sheaf with inverse
$f'/f$. Note that $\psi_{(f')}(f/f') = \psi(f)/s' = s/s'$
and $\psi_{(f)}(f'/f) = \psi(f')/s = s'/s$. We claim there
is a unique ring map
$S_{(ff')} \to \Gamma(Y_{ss'}, \mathcal{O}_Y)$ making the
following diagram commute
$$
\xymatrix{
\Gamma(Y_s, \mathcal{O}_Y) \ar[r] &
\Gamma(Y_{ss'}, \mathcal{O}_Y) &
\Gamma(Y_{s, '} \mathcal{O}_Y) \ar[l]\\
S_{(f)} \ar[r] \ar[u]^{\psi_{(f)}} &
S_{(ff')} \ar[u] &
S_{(f')} \ar[l] \ar[u]^{\psi_{(f')}}
}
$$
It exists because we may use the rule
$x/(ff')^n \mapsto \psi(x)/(ss')^n$, which ``works'' by the formulas
above. Uniqueness follows as $\text{Proj}(S)$ is separated, see
Lemma \ref{lemma-proj-separated} and its proof. This shows that the
morphisms $\varphi_f$ and $\varphi_{f'}$ agree over $Y_s \cap Y_{s'}$.
The restrictions of $\alpha_f$ and $\alpha_{f'}$ agree over
$Y_s \cap Y_{s'}$ because the regular functions $s/s'$ and
$\psi_{(f')}(f/f')$ agree. This proves that the morphisms $\psi_f$
glue to a global morphism from $Y$ into $U_d \subset X$, and
that the maps $\alpha_f$ glue to an isomorphism satisfying the
conditions of the lemma.

\medskip\noindent
We still have to show the pair $(\varphi, \alpha)$ is unique.
Suppose $(\varphi', \alpha')$ is a second such pair.
Let $f \in S_d$. By the commutativity of the diagrams in the lemma we have
that the inverse images of $D_{+}(f)$ under both $\varphi$ and $\varphi'$
are equal to $Y_{\psi(f)}$. Since the opens $D_{+}(f)$ are a basis
for the topology on $X$, and since $X$ is a sober topological
space (see Schemes, Lemma \ref{schemes-lemma-scheme-sober})
this means the maps $\varphi$ and $\varphi'$ are the same
on underlying topological spaces. Let us use $s = \psi(f)$ to
trivialize the invertible sheaf $\mathcal{L}$ over $Y_{\psi(f)}$.
By the commutativity of the diagrams we have that
$\alpha^{\otimes n}(\psi^d_{\varphi}(x)) =
\psi(x) = (\alpha')^{\otimes n}(\psi^d_{\varphi'}(x))$
for all $x \in S_{nd}$. By construction of $\psi^d_{\varphi}$
and $\psi^d_{\varphi'}$ we have
$\psi^d_{\varphi}(x) = \varphi^\sharp(x/f^n) \psi^d_{\varphi}(f^n)$
over $Y_{\psi(f)}$,
and similarly for $\psi^d_{\varphi'}$. By the commutativity of
the diagrams of the lemma we deduce that
$\varphi^\sharp(x/f^n) = (\varphi')^\sharp(x/f^n)$.
This proves that $\varphi$ and $\varphi'$ induce the same morphism
from $Y_{\psi(f)}$ into the affine scheme $D_{+}(f) = \Spec(S_{(f)})$.
Hence $\varphi$ and $\varphi'$ are the same as morphisms.
Finally, it remains to show that the commutativity of the
diagram of the lemma singles out, given $\varphi$, a unique
$\alpha$. We omit the verification.
\end{proof}

\noindent
We continue the discussion from above the lemma.
Let $S$ be a graded ring.
Let $Y$ be a scheme. We will consider {\it triples}
$(d, \mathcal{L}, \psi)$ where
\begin{enumerate}
\item $d \geq 1$ is an integer,
\item $\mathcal{L}$ is an invertible $\mathcal{O}_Y$-module, and
\item $\psi : S^{(d)} \to \Gamma_*(Y, \mathcal{L})$ is a graded
ring homomorphism such that $\mathcal{L}$ is generated by
the global sections $\psi(f)$, with $f \in S_d$.
\end{enumerate}
Given a morphism $h : Y' \to Y$ and a triple
$(d, \mathcal{L}, \psi)$ over $Y$ we can pull it back to the
triple $(d, h^*\mathcal{L}, h^* \circ \psi)$.
Given two triples $(d, \mathcal{L}, \psi)$ and
$(d, \mathcal{L}', \psi')$ with the same integer $d$
we say they are {\it strictly equivalent} if there exists
an isomorphism $\beta : \mathcal{L} \to \mathcal{L}'$
such that $\beta \circ \psi = \psi'$ as graded
ring maps $S^{(d)} \to \Gamma_*(Y, \mathcal{L}')$.

\medskip\noindent
For each integer $d \geq 1$ we define
\begin{eqnarray*}
F_d : \Sch^{opp} & \longrightarrow & \textit{Sets}, \\
Y & \longmapsto &
\{\text{strict equivalence classes of triples }
(d, \mathcal{L}, \psi)
\text{ as above}\}
\end{eqnarray*}
with pullbacks as defined above.

\begin{lemma}
\label{lemma-proj-functor-strict}
Let $S$ be a graded ring.
Let $X = \text{Proj}(S)$.
The open subscheme $U_d \subset X$ (\ref{equation-Ud}) represents the
functor $F_d$ and the triple $(d, \mathcal{O}_{U_d}(d), \psi^d)$
defined above is the universal family (see
Schemes, Section \ref{schemes-section-representable}).
\end{lemma}

\begin{proof}
This is a reformulation of Lemma \ref{lemma-converse-construction}
\end{proof}

\begin{lemma}
\label{lemma-apply}
Let $S$ be a graded ring generated as an $S_0$-algebra by
the elements of $S_1$. In this case the scheme $X = \text{Proj}(S)$
represents the functor which associates to a scheme
$Y$ the set of pairs $(\mathcal{L}, \psi)$, where
\begin{enumerate}
\item $\mathcal{L}$ is an invertible $\mathcal{O}_Y$-module, and
\item $\psi : S \to \Gamma_*(Y, \mathcal{L})$ is a graded
ring homomorphism such that $\mathcal{L}$ is generated by
the global sections $\psi(f)$, with $f \in S_1$
\end{enumerate}
up to strict equivalence as above.
\end{lemma}

\begin{proof}
Under the assumptions of the lemma we have $X = U_1$ and the
lemma is a reformulation of Lemma \ref{lemma-proj-functor-strict} above.
\end{proof}

\noindent
We end this section with a discussion of a functor corresponding
to $\text{Proj}(S)$ for a general graded ring $S$.
We advise the reader to skip the rest of this section.

\medskip\noindent
Fix an arbitrary graded ring $S$. Let $T$ be a scheme.
We will say two triples $(d, \mathcal{L}, \psi)$ and
$(d', \mathcal{L}', \psi')$ over $T$ with possibly different integers
$d$, $d'$ are {\it equivalent} if there exists
an isomorphism
$\beta : \mathcal{L}^{\otimes d'} \to (\mathcal{L}')^{\otimes d}$
of invertible sheaves over $T$
such that $\beta \circ \psi|_{S^{(dd')}}$ and $\psi'|_{S^{(dd')}}$ agree
as graded ring maps $S^{(dd')} \to \Gamma_*(Y, (\mathcal{L}')^{\otimes dd'})$.

\begin{lemma}
\label{lemma-equivalent}
Let $S$ be a graded ring. Set $X = \text{Proj}(S)$. Let $T$ be a scheme.
Let $(d, \mathcal{L}, \psi)$ and $(d', \mathcal{L}', \psi')$
be two triples over $T$. The following are equivalent:
\begin{enumerate}
\item Let $n = \text{lcm}(d, d')$. Write $n = ad = a'd'$. There exists
an isomorphism
$\beta : \mathcal{L}^{\otimes a} \to (\mathcal{L}')^{\otimes a'}$
with the property that
$\beta \circ \psi|_{S^{(n)}}$ and $\psi'|_{S^{(n)}}$ agree
as graded ring maps $S^{(n)} \to \Gamma_*(Y, (\mathcal{L}')^{\otimes n})$.
\item The triples $(d, \mathcal{L}, \psi)$ and $(d', \mathcal{L}', \psi')$
are equivalent.
\item For some positive integer $n = ad = a'd'$ there exists
an isomorphism
$\beta : \mathcal{L}^{\otimes a} \to (\mathcal{L}')^{\otimes a'}$
with the property that
$\beta \circ \psi|_{S^{(n)}}$ and $\psi'|_{S^{(n)}}$ agree
as graded ring maps $S^{(n)} \to \Gamma_*(Y, (\mathcal{L}')^{\otimes n})$.
\item The morphisms $\varphi : T \to X$ and $\varphi' : T \to X$
associated to $(d, \mathcal{L}, \psi)$ and $(d', \mathcal{L}', \psi')$
are equal.
\end{enumerate}
\end{lemma}

\begin{proof}
Clearly (1) implies (2) and (2) implies (3) by restricting to
more divisible degrees and powers of invertible sheaves.
Also (3) implies (4) by the uniqueness statement
in Lemma \ref{lemma-converse-construction}.
Thus we have to prove that (4) implies (1). Assume (4),
in other words $\varphi = \varphi'$.
Note that this implies that we may write
$\mathcal{L} = \varphi^*\mathcal{O}_X(d)$ and
$\mathcal{L}' = \varphi^*\mathcal{O}_X(d')$.
Moreover, via these identifications we have that the graded ring
maps $\psi$ and $\psi'$ correspond to the restriction of the canonical
graded ring map
$$
S \longrightarrow \bigoplus\nolimits_{n \geq 0} \Gamma(X, \mathcal{O}_X(n))
$$
to $S^{(d)}$ and $S^{(d')}$ composed with pullback by $\varphi$
(by Lemma \ref{lemma-converse-construction} again). Hence taking
$\beta$ to be the isomorphism
$$
(\varphi^*\mathcal{O}_X(d))^{\otimes a} =
\varphi^*\mathcal{O}_X(n) =
(\varphi^*\mathcal{O}_X(d'))^{\otimes a'}
$$
works.
\end{proof}

\noindent
Let $S$ be a graded ring.
Let $X = \text{Proj}(S)$.
Over the open subscheme scheme $U_d \subset X = \text{Proj}(S)$
(\ref{equation-Ud}) we have the triple
$(d, \mathcal{O}_{U_d}(d), \psi^d)$. Clearly, if $d | d'$ the triples
$(d, \mathcal{O}_{U_d}(d), \psi^d)$ and
$(d', \mathcal{O}_{U_{d'}}(d'), \psi^{d'})$ are equivalent
when restricted to the open $U_d$ (which is a subset of $U_{d'}$).
This, combined with Lemma \ref{lemma-converse-construction} shows
that morphisms $Y \to X$ correspond roughly to
equivalence classes of triples over $Y$. This is not quite true since if $Y$ is
not quasi-compact, then there may not be a single triple which works.
Thus we have to be slightly careful in defining the corresponding functor.

\medskip\noindent
Here is one possible way to do this. Suppose $d' = ad$.
Consider the transformation of functors $F_d \to F_{d'}$
which assigns to the triple $(d, \mathcal{L}, \psi)$ over
$T$ the triple $(d', \mathcal{L}^{\otimes a}, \psi|_{S^{(d')}})$.
One of the implications of Lemma \ref{lemma-equivalent} is that the
transformation $F_d \to F_{d'}$ is injective!
For a quasi-compact scheme $T$ we define
$$
F(T) = \bigcup\nolimits_{d \in \mathbf{N}} F_d(T)
$$
with transition maps as explained above. This clearly defines a
contravariant functor on the category of quasi-compact schemes
with values in sets. For a general scheme
$T$ we define
$$
F(T)
=
\lim_{V \subset T\text{ quasi-compact open}} F(V).
$$
In other words, an element $\xi$ of $F(T)$ corresponds to a compatible system
of choices of elements $\xi_V \in F(V)$ where $V$ ranges over the
quasi-compact opens of $T$.
We omit the definition of the pullback map $F(T) \to F(T')$
for a morphism $T' \to T$ of schemes.
Thus we have defined our functor
\begin{eqnarray*}
F : \Sch^{opp} & \longrightarrow & \textit{Sets}
\end{eqnarray*}

\begin{lemma}
\label{lemma-proj-functor}
Let $S$ be a graded ring.
Let $X = \text{Proj}(S)$.
The functor $F$ defined above is representable by the scheme $X$.
\end{lemma}

\begin{proof}
We have seen above that the functor $F_d$ corresponds to the
open subscheme $U_d \subset X$. Moreover the transformation
of functors $F_d \to F_{d'}$ (if $d | d'$) defined above
corresponds to the inclusion morphism $U_d \to U_{d'}$
(see discussion above). Hence to show that $F$ is represented
by $X$ it suffices to show that $T \to X$ for a quasi-compact
scheme $T$ ends up in some $U_d$, and that for a general scheme
$T$ we have
$$
\Mor(T, X)
=
\lim_{V \subset T\text{ quasi-compact open}} \Mor(V, X).
$$
These verifications are omitted.
\end{proof}







